\clearpage
\begin{summarybox}
  \textbf{\textcolor{CSummary}{一分钟快记}}\par\nopagebreak[4]

  \textbf{原轴导中，原中导轴；先找中点 $a$，再看函数值或斜率。}

\begin{center}
\begin{tikzpicture}[font=\small,>=stealth,line join=round]
\draw[draw=teal!45!black,fill=green!3,rounded corners=3pt] (0,0) rectangle (13.6,8.2);
\node at (6.8,7.6499999999999995) {在 $a=0$ 的具体模型中记住：原轴导中，原中导轴};
\begin{scope}[shift={(2.8,5.8)},x=1.3cm,y=0.42cm]
\draw[gray,->] (-1.7,0)--(1.8,0) node[right] {$x$};
\draw[gray,->] (0,0)--(0,2);
\draw[teal!75!black,very thick,domain=-1.05:1.05,samples=100] plot(\x,{(\x)^2});
\fill[teal!75!black] (0,0) circle(2.5pt);
\end{scope}
\node at (2.8,4.5) {$x^2$：原函数轴对称};
\begin{scope}[shift={(9.6,5.8)},x=1.3cm,y=0.42cm]
\draw[gray,->] (-1.7,0)--(1.8,0) node[right] {$x$};
\draw[gray,->] (0,-2.1)--(0,2.1);
\draw[blue!70!black,very thick,domain=-1.05:1.05,samples=100] plot(\x,{2*\x});
\fill[blue!70!black] (0,0) circle(2.5pt);
\end{scope}
\node at (9.6,4.5) {$2x$：导函数中心对称};
\begin{scope}[shift={(2.8,2.3)},x=1.3cm,y=0.42cm]
\draw[gray,->] (-1.7,0)--(1.8,0) node[right] {$x$};
\draw[gray,->] (0,-1.2)--(0,1.2);
\draw[teal!75!black,very thick,domain=-1.05:1.05,samples=100] plot(\x,{(\x)^3});
\fill[teal!75!black] (0,0) circle(2.5pt);
\end{scope}
\node at (2.8,0.9999999999999998) {$x^3$：原函数中心对称};
\begin{scope}[shift={(9.6,2.3)},x=1.3cm,y=0.42cm]
\draw[gray,->] (-1.7,0)--(1.8,0) node[right] {$x$};
\draw[gray,->] (0,-0.3)--(0,3.4);
\draw[blue!70!black,very thick,domain=-1.05:1.05,samples=100] plot(\x,{3*(\x)^2});
\fill[blue!70!black] (0,0) circle(2.5pt);
\end{scope}
\node at (9.6,0.9999999999999998) {$3x^2$：导函数轴对称};
\draw[->,teal!65!black,thick] (5.5,6.1)--(6.9,6.1) node[midway,above] {求导};
\draw[->,teal!65!black,thick] (5.5,2.65)--(6.9,2.65) node[midway,above] {求导};
\node at (6.8,.35) {一般位置先找 $a$；中心纵坐标分别看 $0$ 或 $f(a)$。};
\end{tikzpicture}

\end{center}

  \medskip
  \textbf{\textcolor{CSummary}{两条最重要的双向公式}}
  \[
    \boxed{f(a+t)=f(a-t)\ \Longleftrightarrow\ f'(a+t)=-f'(a-t)}
  \]
  \[
    \boxed{f(a+t)+f(a-t)=2f(a)\ \Longleftrightarrow\ f'(a+t)=f'(a-t)}
  \]
  \textbf{对应图像：}原函数关于 $x=a$ 轴对称 $\Longleftrightarrow$ 导函数关于 $(a,0)$ 中心对称；原函数关于 $(a,f(a))$ 中心对称 $\Longleftrightarrow$ 导函数关于 $x=a$ 轴对称。

  \medskip
  \textbf{\textcolor{CSummary}{推广：导函数关于 $(a,b)$ 中心对称}}
  \[
    \boxed{f'(a+t)+f'(a-t)=2b\ \Longleftrightarrow\ f(a+t)-f(a-t)=2bt.}
  \]
  \begin{itemize}[leftmargin=*,itemsep=0.25em,topsep=0.25em]
    \item $b=0$：原函数关于 $x=a$ 对称。
    \item $b\ne0$：\textbf{$x=a$ 不是对称轴}，但可能有其他对称轴。
  \end{itemize}

  \medskip
  \textbf{\textcolor{CSummary}{解题三步}}
  \begin{enumerate}[leftmargin=*,itemsep=0.25em,topsep=0.25em]
    \item 把条件改成 $a+t$ 和 $a-t$ 的配对形式。
    \item 看是左右\textbf{函数值同高／成对}，还是\textbf{导数同号／反号}，选双向公式。
    \item 要求周期时继续连用对称式；只得到一个周期，别急着写最小正周期。
  \end{enumerate}
  \textbf{\textcolor{CSummary}{使用前检查：}}$f$ 在所讨论的同一个对称区间上可导，且每一对对应位置都在该区间中。
\end{summarybox}
